\chapter*{Περίληψη}
\addcontentsline{toc}{chapter}{Περίληψη}

Το παρόν project αποτελεί μια διαδραστική εκπαιδευτική προσομοίωση στο πλαίσιο του προπτυχιακού
μαθήματος ``Νέες Τεχνολογίες στη Διδακτική των ΦΕ'' του Τμήματος Φυσικής του Πανεπιστημίου
Ιωαννίνων. Κύριος στόχος είναι η εξοικείωση των φοιτητών/τριών με τις σύγχρονες συνεργατικές
μεθόδους που χρησιμοποιούνται στην επιστημονική έρευνα και την ανάπτυξη λογισμικού. Μέσα από τη
συνδυαστική χρήση του \LaTeX{} για την επιστημονική συγγραφή, της \texttt{Python (Matplotlib)} για
την προγραμματιστική παραγωγή γραφημάτων, και του \texttt{Git/GitHub} για τον έλεγχο εκδόσεων
(version control) και τη διαχείριση κλάδων (branch management), οι συμμετέχοντες συνεργάζονται
δυναμικά για τη σύνταξη ενός ενιαίου επιστημονικού συγγράμματος Φυσικής.
